% Figure: Hamilton's principle. Among all trial paths between fixed
% endpoints, the physical path makes the action stationary. Trial paths
% are drawn as perturbations of the true path.
\begin{figure}[htbp]
\centering
\begin{tikzpicture}[
  axis/.style={draw=engblack, -{Stealth}, thick},
  truepath/.style={draw=engblue, very thick},
  trial/.style={draw=enggray, thick, dashed},
  endpt/.style={draw=engred, fill=engred, circle, inner sep=1.8pt},
  label/.style={font=\small},
]
% t axis horizontal, q axis vertical
\draw[axis] (0,0) -- (10.4,0) node[anchor=west, label] {$t$};
\draw[axis] (0,0) -- (0,5.4) node[anchor=south, label] {$q$};

% endpoints fixed: (t1, q1) = (1, 1) plotted (1,1); (t2, q2) = (9, 4)
\node[endpt] (P1) at (1,1) {};
\node[endpt] (P2) at (9,4) {};
\node[font=\scriptsize, anchor=north east] at (1,1) {$(t_1,q_1)$};
\node[font=\scriptsize, anchor=north west] at (9,4) {$(t_2,q_2)$};
\draw[draw=enggray, dashed, thin] (1,1) -- (1,0) node[font=\scriptsize, below] {$t_1$};
\draw[draw=enggray, dashed, thin] (9,4) -- (9,0) node[font=\scriptsize, below] {$t_2$};

% true path: smooth curve from (1,1) to (9,4); use a gentle parabola-ish arc
\draw[truepath] plot[domain=1:9, samples=80]
  ({\x}, {1 + 0.375*(\x-1) + 0.6*sin(deg((\x-1)*3.14159/8))});
\node[engblue, label, anchor=south] at (5, 4.4) {physical path (action stationary)};

% trial path 1: above
\draw[trial] plot[domain=1:9, samples=80]
  ({\x}, {1 + 0.375*(\x-1) + 0.6*sin(deg((\x-1)*3.14159/8)) + 0.9*sin(deg((\x-1)*3.14159/8))});
% trial path 2: below
\draw[trial] plot[domain=1:9, samples=80]
  ({\x}, {1 + 0.375*(\x-1) + 0.6*sin(deg((\x-1)*3.14159/8)) - 0.9*sin(deg((\x-1)*3.14159/8))});
\node[font=\scriptsize, enggray, anchor=west] at (5.2, 2.0) {trial paths $q+\delta q$};

% variation arrow at t = 5
\draw[draw=englime, {Stealth}-{Stealth}, thick] (5, {1+0.375*4+0.6+(-0.9)}) -- (5, {1+0.375*4+0.6+0.9});
\node[englime, font=\scriptsize, anchor=west] at (5.05, {1+0.375*4+0.6}) {$\delta q$};

\end{tikzpicture}
\caption[Hamilton's principle and admissible trial paths]{Hamilton's
principle in one coordinate. The endpoints $(t_1,q_1)$ and $(t_2,q_2)$
are fixed (red). Among all admissible trial paths $q(t)+\delta q(t)$
that share those endpoints (grey dashed), the path the system actually
follows (blue) is the one for which the action $S=\int_{t_1}^{t_2}L\,dt$
is stationary: $\delta S=0$ to first order in the variation $\delta q$
(green). The Euler-Lagrange equations are exactly the condition that
this stationarity hold for every admissible variation. Mechanics, cast
this way, is an optimisation problem, and the same variational frame
carries into elasticity, optics, and field theory.}
\label{fig:vol03ch05:action-paths}
\end{figure}
